\clearpage
\section{Reproducibility and Implementation Notes}
\label{app:reproducibility}
\setcounter{equation}{0}\setcounter{table}{0}\setcounter{figure}{0}
\subsection{Environment, records, and reproducibility limits}
The computational archive is the project repository \url{https://github.com/U-Olim/DML-IVQR_Master_Thesis}. This reference identifies the archived workflow; it is not a claim of bit-for-bit regeneration from a seed on arbitrary systems. Preserved source is in \path{simulation/}, \path{Empirical_work/}, and \path{example/}; extension wrappers and configuration are in \path{thesis_extension/src/} and \path{thesis_extension/config/extension_config.R}. The empirical freeze identifies commit \texttt{ac9a9ccad027cde05488abb633e66cc02f87cc46} (2 September 2026).

The canonical environment is R~3.4.3 (30 November 2017) on 64-bit Windows, using five workers. The captured attached packages are \texttt{doSNOW} 1.0.16, \texttt{snow} 0.4-2, \texttt{iterators} 1.0.8, \texttt{foreach} 1.4.3, \texttt{mvtnorm} 1.0-6, \texttt{hqreg} 1.4, \texttt{hdm} 0.2.0, \texttt{quantreg} 5.34, and \texttt{SparseM} 1.77. Environment session and library manifests are under \path{thesis_extension/environment/}. The main simulation uses seed 675; the grid diagnostic also uses seed 675, with saved primitives reloaded in Stage 2; calibration uses seed 20260813. Worker-specific RNG streams are not explicitly initialized in the preserved parallel DML procedure, so the parent seed alone does not establish exact regeneration of every penalty draw.

Frozen CSVs permit deterministic reporting without re-estimation. Input/output manifests in \path{thesis_extension/reporting/}, diagnostic provenance records, reporting consistency checks, and figure validation records connect the reported numbers to archived outputs. Historical manifests describe their archived presentation state and can retain former asset paths. The original figure check covered 480 plotted values; selected-figure validation separately covered 120 coverage values, 120 full-grid frequencies, and 168 median rejection probabilities. The coverage summary is reconstructed from integer direct-at-truth counts; the variance illustration is analytical. This editorial conversion does not rerun those historical validations or simulations.

\subsection{Penalty, density proxy, and residualization}
For each candidate, $b_0$ is the fitted outcome intercept and $b$ the slopes. In the preserved \texttt{rq.fit.lasso} augmented-observation implementation, the unscaled check-loss sum receives a penalty with factor $1/2$. The vector $\lambda^{rq}(a)$ is the supplied routine argument, including an intercept penalty. There is no post-Lasso refit. Let $s_k=(n^{-1}\sum_iX_{ik}^2)^{1/2}$; $U_{ir}$ are independent artificial Uniform$(0,1)$ draws, with $r=1,\ldots,1000$, and $\widehat Q_{0.9}$ is their empirical pivotal-statistic 90th percentile. Fresh draws are made on every penalty call. This is algorithmic randomness, not sample splitting.

In the following formulas, $\phi(e;\mu,\sigma)$ denotes a normal density parameterized by standard deviation. The proxy supplies $\widehat\sigma_\epsilon^2$ in that argument: the preserved R call is \texttt{dnorm(e, mean(e), var(e))}. It evaluates a fitted-normal density at residuals rather than estimating the theoretical conditional density at zero. The instrument penalty is $\mathcal P_{j,a}(\gamma_j)=\sum_k\lambda^\Gamma_{jk}(a)|\gamma_{jk}|$, using default iterative heteroskedastic loadings. Its ordinary intercept is unpenalized; no post-Lasso refit is performed.
\begin{equation}
    \begin{pmatrix}
        \widehat b_0(a)\\
        \widehat\beta(a)
    \end{pmatrix}
    =
    \operatorname*{arg\,min}_{b_0,b}
    \left\{
        \sum_{i=1}^{n}
        \rho_{\tau}\!\left(
            Y_i-D_i^{\prime}a-b_0-X_i^{\prime}b
        \right)
        +
        \frac{1}{2}\left[
            \lambda_0^{rq}(a)|b_0|
            +\sum_{k=1}^{p}\lambda_k^{rq}(a)|b_k|
        \right]
    \right\}.
    \label{eq:beta-hat-dml}
\end{equation}

\begin{equation}
    \begin{aligned}
        T_r
        &=
        \max_{1\leq k\leq p}
        \left|
            \frac{
                \sum_{i=1}^{n}X_{ik}
                \left[\tau-\mathbf{1}\{U_{ir}<\tau\}\right]
            }{
                s_k\sqrt{\tau(1-\tau)}
            }
        \right|,\\
        \lambda^{rq}(a)
        &=
        2\,\widehat Q_{0.9}\!\left(\{T_r\}_{r=1}^{1000}\right)
        \sqrt{\tau(1-\tau)}
        (1,s_1,\ldots,s_p)^{\prime}.
    \end{aligned}
    \label{eq:bc-pivotal-penalty}
\end{equation}

\begin{equation}
    \begin{aligned}
        \widehat{\epsilon}_i(a)
        &=
        Y_i-D_i^{\prime}a-\widehat b_0(a)
        -X_i^{\prime}\widehat{\beta}(a),\\
        \widehat{\mu}_{\epsilon}(a)
        &=\frac{1}{n}\sum_{i=1}^{n}\widehat{\epsilon}_i(a),\\
        \widehat f_i(a)
        &=
        \phi\!\left(
            \widehat{\epsilon}_i(a);
            \widehat{\mu}_{\epsilon}(a),
            \widehat{\sigma}_{\epsilon}^{2}(a)
        \right),
    \end{aligned}
    \label{eq:estimated-profile-residual}
\end{equation}

\begin{equation}
    \begin{pmatrix}
        \widehat r_{j0}(a)\\
        \widehat\Gamma_{\mathrm{imp},j}(a)
    \end{pmatrix}
    =
    \operatorname*{arg\,min}_{(r_{j0},\gamma_j)}
    \left\{
        \sum_{i=1}^{n}
        \left[
            \sqrt{\widehat f_i(a)}Z_{ji}
            -r_{j0}
            -\sqrt{\widehat f_i(a)}X_i^{\prime}\gamma_j
        \right]^2
        +\mathcal{P}_{j,a}(\gamma_j)
    \right\}.
    \label{eq:dml-residualization-implementation}
\end{equation}

\begin{equation}
    \begin{aligned}
        \widehat\Psi_{ji}(a)
        &=
        Z_{ji}-\widehat r_{j0}(a)-X_i^{\prime}\widehat\Gamma_{\mathrm{imp},j}(a),\\
        \widehat\Psi_i(a)
        &=
        Z_i-\widehat r_0(a)-\widehat\Gamma_{\mathrm{imp}}(a)X_i.
    \end{aligned}
    \label{eq:estimated-residualized-instrument}
\end{equation}

\begin{equation}
    \begin{pmatrix}
        \widehat b_{0,\mathrm{orc}}(a)\\
        \widehat\beta_{\mathrm{orc}}(a)
    \end{pmatrix}
    =
    \operatorname*{arg\,min}_{b_0,b}
    \sum_{i=1}^{n}
    \rho_{\tau}\!\left(
        Y_i-D_i^{\prime}a-b_0-X_{S_0,i}^{\prime}b
    \right).
    \label{eq:oracle-beta}
\end{equation}

\begin{equation}
    \begin{pmatrix}
        \widehat b_{0,\mathrm{full}}(a)\\
        \widehat\beta_{\mathrm{full}}(a)
    \end{pmatrix}
    =
    \operatorname*{arg\,min}_{b_0,b}
    \sum_{i=1}^{n}
    \rho_{\tau}\!\left(
        Y_i-D_i^{\prime}a-b_0-X_i^{\prime}b
    \right).
    \label{eq:full-beta}
\end{equation}

\begin{equation}
    \begin{aligned}
        \widehat M_e(a)
        &=\frac{1}{n}\sum_{i=1}^{n}
          Z_iX_{e,i}^{\prime}\widehat f_{e,i}(a),&
        \widehat J_e(a)
        &=\frac{1}{n}\sum_{i=1}^{n}
          X_{e,i}X_{e,i}^{\prime}\widehat f_{e,i}(a),\\
        \widehat\delta_e(a)
        &=\widehat M_e(a)\widehat J_e(a)^{-1},&
        \widehat\Psi_{e,i}(a)
        &=Z_i-\widehat\delta_e(a)X_{e,i}.
    \end{aligned}
    \label{eq:benchmark-residualization}
\end{equation}
The first five equations describe DML-IVQR. The last three describe Oracle-GMM and Full-GMM, with $S_0=\{1,\ldots,10\}$, $X_{\mathrm{orc},i}=X_{S_0,i}$, and $X_{\mathrm{full},i}=X_i$. Their quantile regressions include an intercept, but their matrix residualization includes no constant regressor. DML retains the ordinary intercept of a transformed regression and subtracts it from the original instrument without dividing by $\sqrt{\widehat f_i(a)}$. It is therefore not the exact sample counterpart of the ideal density-weighted intercept. All procedures use same-sample nuisance fits and scores; cross-fitting is absent. The comparison does not isolate regularization from these residualization differences.

\subsection{Score algebra and its conditions}
\label{app:score-algebra}
The ideal nuisance vector and derivative identities are retained below. They presume finite density-weighted moments and permissible differentiation, conditions not established for the simulation DGP's non-finite true residual density. They are not implementation-level assertions.
\begin{equation}
    \eta(a)
    =
    \begin{pmatrix}
        c(a)\\
        \beta(a)\\
        \operatorname{vec}\{\delta(a)\}
    \end{pmatrix},
    \qquad
    \eta_0
    =
    \eta\!\left(\alpha_0(\tau)\right)
    =
    \begin{pmatrix}
        c_0(\tau)\\
        \beta_0(\tau)\\
        \operatorname{vec}\{\delta(\alpha_0(\tau))\}
    \end{pmatrix}.
    \label{eq:profile-eta}
\end{equation}

\begin{equation}
    \operatorname{E}\!\left[
        \left\{
            \tau-\mathbf{1}\!\left[
                Y_i\leq D_i^{\prime}a+c(a)+X_i^{\prime}\beta(a)
            \right]
        \right\}\widetilde X_i
    \right]
    =0.
    \label{eq:profile-beta-moment}
\end{equation}

\begin{equation}
    \begin{aligned}
        M_0
        &=\operatorname{E}\!\left[
            Z_i\widetilde X_i^{\prime}
            f_{\epsilon_0}(0\mid X_i,Z_i)
        \right],\\
        J_0
        &=\operatorname{E}\!\left[
            \widetilde X_i\widetilde X_i^{\prime}
            f_{\epsilon_0}(0\mid X_i,Z_i)
        \right],\\
        \delta_0&=M_0J_0^{-1}.
    \end{aligned}
    \label{eq:orthogonality-MJ}
\end{equation}

\begin{equation}
    \begin{aligned}
        \left.
        \partial_{(c,\beta^{\prime})}
        \operatorname{E}\!\left[
            g_{\tau}\!\left(W_i;\alpha_0(\tau),\eta\right)
        \right]
        \right|_{\eta=\eta_0}
        &=-\operatorname{E}\!\left[
            (Z_i-\delta_0\widetilde X_i)\widetilde X_i^{\prime}
            f_{\epsilon_0}(0\mid X_i,Z_i)
        \right]\\
        &=-\left[M_0-\delta_0J_0\right]
        =0.
    \end{aligned}
    \label{eq:orthogonality-beta}
\end{equation}

\begin{equation}
    \left.
    \partial_{\delta_j}
    \operatorname{E}\!\left[
        g_{\tau,j}\!\left(W_i;\alpha_0(\tau),\eta\right)
    \right]
    \right|_{\eta=\eta_0}
    =-\operatorname{E}\!\left[
        \left\{
            \tau-\mathbf{1}\!\left[
                Y_i\leq D_i^{\prime}\alpha_0(\tau)
                +c_0(\tau)
                +X_i^{\prime}\beta_0(\tau)
            \right]
        \right\}\widetilde X_i^{\prime}
    \right]
    =0.
    \label{eq:orthogonality-delta}
\end{equation}
The first derivative cancels by $\delta_0=M_0J_0^{-1}$; the second cancels by the profiled quantile moment. Orthogonality concerns local nuisance sensitivity, and does not establish identification or finite-sample calibration.
